% Generated draft: requires mathematical and source review.
\documentclass[11pt]{article}
\usepackage[margin=25mm]{geometry}
\usepackage{fontspec}
\setmainfont{Latin Modern Roman}
\usepackage{amsmath,amssymb,mathtools,unicode-math}
\setmathfont{Latin Modern Math}
\usepackage{enumitem,needspace,xurl,hyperref}
\setlength{\parindent}{0pt}
\setlength{\parskip}{5pt}
\setlength{\emergencystretch}{3em}
\newcommand{\N}{\mathbb{N}}
\newcommand{\Z}{\mathbb{Z}}
\newcommand{\Q}{\mathbb{Q}}
\newcommand{\R}{\mathbb{R}}
\newcommand{\C}{\mathbb{C}}
\newcommand{\F}{\mathbb{F}}
\newcommand{\A}{\mathbb{A}}
\newcommand{\PP}{\mathbb{P}}
\newcommand{\E}{\mathbb{E}}
\newcommand{\eps}{\varepsilon}
\newcommand{\dd}{\,\mathrm{d}}
\begin{document}
\section*{UnsolvedMath ID: 20000222}
\textbf{Source problem code:} catalogue entry
\subsection{Definitions and mathematical statement}
The damaged AIM extraction asks whether $X$ has a linear series $g^r_{d_1,d_2}$ exactly when $\rho(g,r,d_1+d_2)\geq0$, subject to a suitable ``relatively balanced'' condition. In the binary-curve interpretation supplied by the research summary, $X$ is a genus-$g$ curve obtained by joining two copies of $\PP^1$ at $g+1$ paired distinct nodes. A multidegree-$(d_1,d_2)$ line bundle has those degrees on the two components, and a $g^r_{d_1,d_2}$ is a pair $(L,V)$ with $\dim V=r+1$ and $V\subset H^0(X,L)$. The original extraction does not supply the preceding definition of $X$, a generality assumption, or base-point-freeness conventions, so this interpretation requires confirmation.

The Brill--Noether number is
\[
 \rho(g,r,d)=g-(r+1)(g-d+r),\qquad d=d_1+d_2.
\]
The source discusses Caporaso's weaker balance condition and alternatively the explicit bounds $0\leq d_i\leq g-1$ for $i=1,2$. It expressly allows different answers for different balance definitions. The weakest condition itself is not defined in the supplied text. The intended existence equivalence therefore remains dependent on both the choice of balance and the precise general binary-curve setting.

The supplied report claims counterexamples to sufficiency of both proposed balance conditions: for every integer $r\geq1$, take
\[
 g=2(r+1),\qquad(d_1,d_2)=(r-1,2r+1),\qquad d=3r.
\]
It states that both balance requirements hold and $\rho=0$, but no $g^r_{d_1,d_2}$ exists. This is a negative claim about those criteria, not a proof that no refined balance criterion can work.

Its positive claim assumes an algebraically closed field of characteristic zero. For integers $r,e\geq1$, put $m=re$. On a general genus-$g$ binary curve, whenever $g\leq(r+1)e$, it claims a nonempty rational open locus of $g^r_{r,m}$'s isomorphic to an open subset of
\[
 \PP^{\rho},\qquad \rho=\rho(g,r,r+re)=r\bigl((r+1)e-g\bigr).
\]
At the boundary $g=(r+1)e$, hence $\rho=0$, the report claims that the full fixed-multidegree set consists of one isomorphism class, set-theoretically. This is not a scheme-theoretic reducedness, smoothness or multiplicity-one assertion, nor a classification for arbitrary multidegrees.

The source motivates the question by degenerations for Brill--Noether generality and by dimensions of limit linear series on non-compact-type and pseudocompact curves. Its trailing Remark 4.2 is incomplete (``This is closely related to'') and contributes no recoverable additional statement.
\subsection{Short English statement}
Ask whether Brill--Noether nonnegativity characterizes linear-series existence on the intended general binary curve for an appropriately balanced multidegree. The report gives a counterexample family to two naive balance choices and a positive rational interpolation locus for multidegrees $(r,re)$, with a unique set-theoretic class at its zero-dimensional boundary.
\subsection{Sources}
\begin{itemize}
\item Local registry: \url{data/unsolvedmath/problems.json}, permanent ID $20000222$, source code \texttt{catalogue entry}; damaged original extraction, incomplete remark and full research summary. Statement status is reconstructed and unverified.
\item AIM workshop ``Degenerations in algebraic geometry'', Problem 4.1: \url{https://aimath.org/pastworkshops/degenalggeomproblems.pdf}. Canonical location \texttt{aim-algebraic-geometry-notes.json}, \texttt{notes[221]}. Supplied verification says the PDF recovers notation and a preceding definition of $X$, but that preceding text is not included in this record; no PDF was opened here.
\item Original text cites Caporaso, ``Brill-Noether theory of binary curves'', and attributes results for $r\leq2$ using a weak balance notion. No further bibliographic data or exact theorem supplied; original literature field empty.
\item Supplied machine-generated research attempt 1: partial result; problem partially solved at the run and in the registry; rigor metadata ``proved'', novelty confidence low. This non-peer-reviewed research aid requires independent verification before citation.
\end{itemize}
\end{document}
